
As AI systems improve on human preference benchmarks, do the humans interacting with them become less engaged --- submitting shorter prompts, correcting errors less often, and discriminating less carefully among model outputs? This Bidirectional Alignment Asymmetry (BAA) conjecture predicts declining human behavioral engagement as a consequence of AI quality improvement, but it has not been tested empirically at scale. We construct a returning-user cohort from WildChat-1M (27,902 users defined by IP-hash pseudonymization, 13 monthly bins, April 2023--April 2024) and apply autocorrelation-corrected Mann-Kendall trend tests to three behavioral proxies: prompt token count, preference vote entropy, and correction frequency. Contrary to the BAA directional prediction for prompt length --- and noting that the causal chain from AI quality improvement to behavioral change is not verified in this study --- prompt token count increases strongly over the observation period ($\tau$ = +0.744, 95\% CI [0.415, 0.972], p = 0.0005, Hamed-Rao autocorrelation-corrected, ACF lag-1 = 0.634): returning users compose progressively longer prompts across 13 monthly bins. The remaining two proxies are unmeasurable under current conditions: the LMSYS Chatbot Arena primary dataset (\texttt{lmsys/chatbot\textbackslash \_arena\textbackslash \_conversations}) is access-gated on HuggingFace, and the fallback dataset contains no timestamps, making temporal vote entropy analysis impossible; explicit correction behavior occurs in approximately 0.05\% of turns, too sparse for trend detection. The BAA disengagement directional prediction for prompt length is not supported in this cohort. The result is a directional refutation within a self-selected high-engagement population; the causal question of whether AI quality improvement drives behavioral change remains empirically unresolved. We document the measurement pipeline and infrastructure constraints as a foundation for future studies.

